\chapter{Representations of quantum mechanics}\label{ann:pictures}
\section*{Brief complement of quantum mechanics}
An average is defined in quantum mechanics by
\begin{align*}
    \expval{O}_\psi&=\expval{O}{\psi},\\
    \expval{O}_\rho&=\Tr{\rho O},
\end{align*}
where the subscripts $\ket{\psi}$ and $\rho$ refer respectively to the representation of pure states or
to density matrix more general formalism. 
\section{Schrödinger and Heisenberg pictures}\noindent
In the Schrödinger picture we saw that $\ket{\psi(t)}=U(t)\ket{\psi}$, where $U(t)=e^{-\frac{i}{\hbar}Ht}$ in the case of a time-independent Hamiltonian with $t_0=0$. Nevertheless, the following results can be derived in the general case of a time dependent Hamiltonian,~\cite{Breuer}. Considering the average value of an observable $A_S(t)$ in the Schrödinger picture (in this picture, the time-dependence of an observable can only be explicit), we find
\begin{align}
    \expval{A(t)} &= \expval{A_S(t)}{\psi(t)}\nonumber\\
    &= \expval{U^\dagger(t,t_0)A_S(t)U(t,t_0)}{\psi(t_0)}\\
    &= \expval{A_H(t)}{\psi}_H.
\end{align}
This calculation naturally reveals the Heisenberg representation of observables and states. Indeed, we have that $\ket{\psi(t_0)}_S\equiv\ket{\psi}_H$ which links both pictures at the initial time $t_0$. Furthermore, we can link both representation at all time by calculating the evolution of an observable in the Heisenberg picture:
\begin{equation}
    \boxed{\dert{A_H} = \frac{i}{\hbar}\left[H_{H},A_{H}\right] + \derpart{A_S}\label{eq:HSlink}}.
\end{equation}
This equation leads to the famous \textit{Ehrenfest theorem} analogous in quantum mechanics:
\begin{equation}
    \boxed{\dert{\expval{A}}=\frac{i}{\hbar}\expval{\qty[H(t),A(t)]}+\expval{\derpart{A}}}.
\end{equation}
We observe that if the observable doesn't depend explicitly on time, then the dynamics of the operator is fully described in the Heisenberg picture.

\section{Interaction picture}\noindent
As mentioned in the Section~\ref{OQS}, we can represent the coupling between $S$ and $E$ with an interaction Hamiltonian such that
\begin{align}
    H_T = \overbrace{H_S\otimes\mathbb{I}_E+\mathbb{I}_S\otimes H_E}^{H_0}+H_I.
\end{align}
Knowing that $\ket{\psi(t)}_S=U(t,t_0)\ket{\psi}_S$, we can write that
\begin{equation}\label{eq:evolrhoU}
    \rho_S(t) = U(t,t_0)\rho_{S}(t_0)U^\dagger(t,t_0).
\end{equation}
The evolution operator $U$ is the solution of the Schrödinger equation for the total system $T$. We know that the free system (dynamically governed by $H_0$) also follows unitary dynamics, which is described by $U_0$. Thus, we can define the interaction unitary operator $U_I$ as
\begin{align}
    U_I(t,t_0) &= U_0^\dagger(t,t_0)U(t,t_0)\\
    \Leftrightarrow U(t,t_0) &= U_0(t,t_0)U_I(t,t_0)\label{eq:Utotal}
\end{align}
If we now try to calculate the average of an operator (in the total space $\mathcal{H}_T$) using the Schrödinger picture, we can find a new interesting relation by injecting the Eq.~(\ref{eq:Utotal}) in the Eq.~(\ref{eq:evolrhoU}) and then into the following average:
\begin{align*}
    \expval{O(t)} &= \Tr{\rho_{S}(t)O_S(t)}\\
    &= \mathrm{Tr}\,[\underbrace{U_I(t,t_0)\rho_{S}(t_0)U_I^\dagger(t,t_0)}_{\tilde{\rho}(t)}\,\underbrace{U_0^\dagger(t,t_0)O_S(t)U_0(t,t_0)}_{\tilde{O}(t)}].
\end{align*}
We can thus define in the interaction picture:
\begin{align}
    \tilde{\rho}(t)&=U_I(t,t_0)\rho_{S}(t_0)U_I^\dagger(t,t_0),\label{eq:rhoIS}\\
    \tilde{O}(t) &= U_0^\dagger(t,t_0)O_S(t)U_0(t,t_0)\label{eq:obsIS}.
\end{align}
We can now reverse the Eq.~(\ref{eq:evolrhoU}) and then use the Eq.~(\ref{eq:Utotal}) to get
\begin{equation*}
    \rho_S(t_0) = U^\dagger(t,t_0)\rho_S(t)U(t,t_0) = U_I^\dagger(t,t_0)U^\dagger_0(t,t_0)\rho_S(t)U_0(t,t_0)U_I(t,t_0),
\end{equation*}
and, by injecting this into the Eq.~(\ref{eq:rhoIS}), we get
\begin{equation}
    \tilde{\rho}(t) = U^\dagger_0(t,t_0)\rho_S(t)U_0(t,t_0).
\end{equation}
This result is obvious because a density matrix is also an operator and thus respects the Eq.~(\ref{eq:obsIS}).
In the end, we find the results used in the main part of this document:
\begin{equation}
    \boxed{\tilde{\rho}(t)=U_I(t,t_0)\rho_{S}(t_0)U_I^\dagger(t,t_0)= U^\dagger_0(t,t_0)\rho_S(t)U_0(t,t_0)}.
\end{equation}
By working on the Liouville-von Neumann equation, one can show that by jumping into the interaction picture the dynamics do not depend anymore on the free Hamiltonian $H_0$:
\begin{equation}
    \qty(\dert{\rho}=-\frac{i}{\hbar}\qty[H_T,\rho])\Leftrightarrow\qty(\dert{\tilde{\rho}}=-\frac{i}{\hbar}\qty[\tilde{H}_I,\tilde{\rho}]),
\end{equation}
where $\tilde{H}_I=U_0^\dagger(t,t_0)H_I U_0(t,t_0)$ as defined. It is thus usual to work in this representation to get rid of free dynamics terms, and then jump back to the Schrödinger picture to recover the proper dynamics of the system.
\subsubsection{Further discussion}\noindent
As seen in Eq.~(\ref{eq:rhoIS}), a state evolves in the interaction picture with a relation depending on $U_I$, which has been defined as $U_I=U_0^\dagger U$. It is possible to show from this definition that $U_I$ is subject to the following differential equation
\begin{equation}
    \dert{U_I}=-\frac{i}{\hbar}\tilde{H}_I U_I,
\end{equation}
where we recall that $\tilde{H}_I=U_0^\dagger H_I U_0$ by definition of the interaction picture. If we now consider the evolution  